\documentclass[a4paper,11pt]{article}
\usepackage[utf8]{inputenc}
\usepackage[T1]{fontenc}
\usepackage[french]{babel}
\usepackage[left=1cm,right=1cm,top=.5cm,bottom=0cm]{geometry}
\usepackage{enumitem}
\usepackage{hyperref}
\usepackage{xcolor}
\usepackage{titlesec}
\usepackage{fontawesome5}
\usepackage{lmodern}
\usepackage[normalem]{ulem}

% --- Couleurs Thales ---
\definecolor{primary}{RGB}{35, 55, 123}
\definecolor{secondary}{RGB}{80, 80, 80}

% --- Configuration des liens ---
\hypersetup{
    colorlinks=true,
    linkcolor=primary,
    filecolor=primary,
    urlcolor=primary,
}

% --- Formatage des titres ---
\titleformat{\section}{\large\bfseries\color{primary}\uppercase}{}{0em}{}[\titlerule]
\titlespacing{\section}{0pt}{8pt}{5pt}

% --- Commandes personnalisées ---
\newcommand{\resumeItem}[1]{\item\small{#1}}
\newcommand{\resumeSubheading}[4]{
  \vspace{-1pt}\item
    \begin{tabular*}{0.97\textwidth}[t]{l@{\extracolsep{\fill}}r}
      \textbf{#1} & \textit{\small #2} \\
      \textit{\small #3} & \textit{\small #4} \\
    \end{tabular*}\vspace{-5pt}
}

\begin{document}

% --- EN-TÊTE ---
\begin{center}
    {\Large \textbf{Loïc Aron MBASSI EWOLO}} \\ \vspace{4pt}
    \textbf{\large Stage : Ingénieur développement FPGA (MPSoC / RFSoC)} \\
    \textit{\small Disponible dès \textbf{Avril 2026} pour 4 à 5 mois (Stage Assistant Ingénieur)}\\
    \vspace{4pt}
    \small
    \faMapMarker\ Cergy, France (Mobile Laval) \quad
    \faPhone\ (+33) 7 59 52 33 98 \quad
    \faEnvelope\ \href{mailto:loicaron.mbassiewolo@ensea.fr}{loicaron.mbassiewolo@ensea.fr} \\ \vspace{2pt}
    \faLinkedin\ \href{https://www.linkedin.com/in/loic-mbassi}{loic-mbassi} \quad
    \faGithub\ \href{https://github.com/Nameless0l}{github} \quad
    \faGlobe\ \href{https://mbassiloic.tech/\#portfolio}{Portfolio : mbassiloic.tech}
\end{center}

\vspace{-6pt}

% --- PROFIL ---
\noindent\small{Élève-ingénieur en double diplôme \textbf{ENSEA/Polytechnique Yaoundé}, spécialisé en \textbf{systèmes embarqués} et \textbf{conception numérique FPGA}. Expérience pratique en développement \textbf{VHDL} sur Cyclone V avec intégration de soft-processeurs (\textbf{Nios V/SoPC}), protocoles I2C/SPI et validation sur cible. Je recherche un stage chez \textbf{Thales} pour contribuer au développement de bancs de tests FPGA de nouvelle génération (RFSoC/MPSoC).}

% --- FORMATION ---
\section{Formation}
\begin{itemize}[leftmargin=*, label={.}, nosep]
    \item \textbf{ENSEA (École Nationale Supérieure de l'Électronique et de ses Applications)} \hfill \textit{Cergy, France | 2025 -- Présent} \\
    \small{Ingénieur 2A, Double Diplôme - \textbf{Systèmes Embarqués}, Traitement du Signal et IA.} \\
    \textit{\small{Projets actuels : Conception numérique FPGA (VHDL, Quartus), Intégration Nios V, Traitement du Signal.}}
    \vspace{2pt}
    \item \textbf{École Nationale Supérieure Polytechnique} \hfill \textit{Yaoundé, Cameroun | 2021 -- 2025} \\
    \small{Diplôme d'Ingénieur de Conception (Niveau Master 2) : Mathématiques Appliquées, Génie Logiciel \& Systèmes.}
    \item \textbf{Université de Yaoundé I} \hfill \textit{Yaoundé, Cameroun | 2020 -- 2022} \\
    \small{Licence en Informatique et Mathématiques : Algorithmique C, Analyse, Algèbre Linéaire.}
\end{itemize}

% --- COMPÉTENCES TECHNIQUES ---
\section{Compétences Techniques}
\begin{itemize}[leftmargin=*, nosep]
    \item \small{\textbf{FPGA \& HDL :} \textbf{VHDL}, Quartus Prime, \textbf{Platform Designer (SoPC)}, Nios V, ModelSim, Cyclone V.}
    \item \small{\textbf{Embarqué \& Hardware :} \textbf{C / C++}, STM32, Protocoles (\textbf{I2C}, SPI, JTAG), Linux Embarqué, Soudure.}
    \item \small{\textbf{Maths \& Algorithmes :} Algèbre Linéaire, Traitement du Signal, SLAM, Optimisation.}
    \item \small{\textbf{IA \& Data :} PyTorch, TensorFlow, Computer Vision (YOLOv10), TinyML.}
    \item \small{\textbf{Outils :} Python, Git, Docker, CI/CD, Altium/KiCad (Conception PCB).}
\end{itemize}

% --- PROJETS TECHNIQUES FPGA ---
\section{Projets FPGA \& Systèmes Embarqués}
\begin{itemize}[leftmargin=*, label={}]

\item \textbf{Système SoPC sur Cyclone V -- Intégration Nios V et Périphériques} \hfill \textit{TP ENSEA (En cours)}
\vspace{-2pt}
\item \small{Conception complète d'un système sur puce programmable avec soft-processeur et périphériques I2C.}
    \begin{itemize}[leftmargin=0.4cm, nosep, label=\tiny$\bullet$]
        \item \small{\textbf{Architecture SoPC :} Instanciation du processeur \textbf{Nios V} via \textbf{Platform Designer}, configuration mémoire On-Chip (128Ko), mapping d'adresses et interconnexion des bus (Avalon).}
        \item \small{\textbf{Développement VHDL :} Description RTL du top-level, gestion des signaux I/O (PIO), contraintes de placement/routage (QSF), simulation comportementale.}
        \item \small{\textbf{Intégration I2C :} Contrôleur Avalon I2C Master pour communication avec accéléromètre ADXL345, gestion des signaux open-drain (SDA/SCL).}
        \item \small{\textbf{Programmation C Embarquée :} Développement BSP/HAL, drivers capteurs, debug JTAG UART.}
    \end{itemize}
    \vspace{3pt}

\item \textbf{Niveau à Bulles Numérique -- Acquisition Capteur et Affichage Temps Réel} \hfill \textit{Projet ENSEA}
\vspace{-2pt}
\item \small{Application embarquée exploitant l'accéléromètre via le soft-processeur Nios V.}
    \begin{itemize}[leftmargin=0.4cm, nosep, label=\tiny$\bullet$]
        \item \small{Lecture registres ADXL345 via API I2C, conversion données 16 bits signés.}
        \item \small{Calcul d'inclinaison et mapping sur LEDs via PIO (\textbf{IOWR\_ALTERA\_AVALON\_PIO\_DATA}).}
        \item \small{Validation fonctionnelle sur carte DE10-Nano (Cyclone V SoC).}
    \end{itemize}
    \vspace{3pt}

\item \textbf{Écran Magique -- Système de Contrôle par Mouvement} \hfill \textit{Projet ENSEA}
\vspace{-2pt}
\item \small{Extension du projet SoPC pour détection de retournement et effacement d'écran.}
    \begin{itemize}[leftmargin=0.4cm, nosep, label=\tiny$\bullet$]
        \item \small{Ajout de PIO supplémentaire pour signal d'effacement, modification du QSYS et régénération BSP.}
        \item \small{Algorithme de détection de seuil sur axe Z de l'accéléromètre.}
    \end{itemize}
    \vspace{3pt}

\item \textbf{Harmony Gloves -- Traduction de la Langue des Signes} \hfill \textit{Laboratoire IOTA (Polytechnique)}
\vspace{-2pt}
\item \small{Prototypage d'un gant instrumenté avec capteurs IMU pour reconnaissance gestuelle.}
    \begin{itemize}[leftmargin=0.4cm, nosep, label=\tiny$\bullet$]
        \item \small{\textbf{Électronique :} Conception circuit, \textbf{soudure} composants (Capteurs Flex/IMU), interfaçage STM32.}
        \item \small{\textbf{Embarqué :} Acquisition données capteurs via I2C/SPI, fusion de données, modèle TinyML.}
    \end{itemize}
    \vspace{3pt}

\item \textbf{Upgrade Jetson -- Robotique Autonome (Challenge Défense)} \hfill \textit{Challenge CoHoMa}
\vspace{-2pt}
\item \small{Programmation d'un robot autonome pour cartographie en environnement non structuré.}
    \begin{itemize}[leftmargin=0.4cm, nosep, label=\tiny$\bullet$]
        \item \small{Implémentation d'algorithmes \textbf{SLAM} et méthodes de \textbf{Monte Carlo} sur Nvidia Jetson.}
        \item \small{Optimisation IA embarquée (YOLOv10), containerisation Docker.}
    \end{itemize}
    \vspace{3pt}

\item \textbf{Projet Taxi -- Simulation Système Temps Réel en C} \hfill \textit{Projet Académique}
\vspace{-2pt}
\item \small{Développement d'un noyau de simulation multiprocessus pour flotte de véhicules.}
    \begin{itemize}[leftmargin=0.4cm, nosep, label=\tiny$\bullet$]
        \item \small{\textbf{Programmation Bas Niveau :} Mémoire partagée (SHM), IPC, synchronisation Signaux UNIX.}
        \item \small{Ordonnanceur FIFO avec prévention des interblocages, visualisation OpenGL.}
    \end{itemize}

\end{itemize}

% --- EXPÉRIENCE PROFESSIONNELLE & RECHERCHE ---
\section{Expériences Professionnelles}
\begin{itemize}[leftmargin=*, label={}]

    \resumeSubheading
      {Assistant de Recherche -- Généralisation IA (Grokking)}{Juin 2025 -- Sept. 2025}
      {MILA (Institut Québécois d'IA) - Collaboration à distance}{\textbf{Québec, Canada}}
      \begin{itemize}[leftmargin=0.4cm, nosep, label=\tiny$\bullet$]
        \item \small{Implémentation de pipelines de tests statistiques sous \textbf{PyTorch}.}
        \item \small{Reproduction d'expériences SOTA et analyse mathématique de la convergence.}
      \end{itemize}
    \vspace{2pt}

    \resumeSubheading
      {Ingénieur Backend (Freelance)}{Oct. 2024 -- Août 2025}
      {CSC Fintech -- Architecture Bancaire Sécurisée}{Douala, Cameroun}
      \begin{itemize}[leftmargin=0.4cm, nosep, label=\tiny$\bullet$]
        \item \small{Architecture Backend robuste (Docker, Redis), gestion des accès concurrents.}
        \item \small{Tests unitaires et revues de code systématiques.}
      \end{itemize}
\end{itemize}

% --- ENGAGEMENT ASSOCIATIF & LEADERSHIP ---
\section{Engagement Associatif \& Certifications}
\begin{itemize}[leftmargin=*, nosep]
    \item \textbf{Leadership :} \textbf{Président du Club Informatique}, \textbf{Chef de la Cellule Projet IA} (Polytechnique Yaoundé).
    \item \textbf{Hackathons :} \textbf{2\textsuperscript{e} Place} - IndabaX Africa (IA Santé), \textbf{1\textsuperscript{er} National} - Mountain Hub.
    \item \textbf{Certifications :} TinyML (IA Frugale - En cours), Supervised ML (DeepLearning.AI), Python (IBM).
    \item \textbf{Langues :} Français (Natif), Anglais (Professionnel technique - Documentation, articles).
\end{itemize}

\end{document}
